\section{Introduction}
\subsection{Identifying \LaTeX's key accessibility problems for screen reader users}
Similar to a tool like Microsoft Word, \latexspaced is a tool used to process documents. \latexspaced does this in a programmatical manner, likening it to a programming language, whereas Microsoft Word has a \acrfull{gui} which simplifies the process for an ordinary user. However, \latexspaced prevails in its ability to process scientific documents which may require graphs, images, tables, or advanced mathematical capabilities.

\subsection{Why accessibility is important}


\parencite{accessibility_importance}

- who would benefit from it
        - stats with disabilities e.g. how many people are blind and use screen readers
            - (use good sources e.g. government websites or world health organisations)
        - severity and frequency are considerations
\subsection{Why LuaLaTeX provides better accessibility than PDFLaTeX}
\pdflatexspaced was created to be a media printer, such that it produced PDF documents with the sole purpose of being printed out. Due to its innate design, this means it lacked support for screen reader users. Additionally, as its focus has shifted to maintenance rather than development, \pdflatexspaced is unsuitable for accessibility where requirements change frequently, making it an outdated \texspaced engine.

The introduction of \xetexspaced improved font handling in \latex , yet its lifetime was short-lived, with it being considered an evolutionary deadend, at the addition of \lualatex.

Unlike \pdflatex, \lualatexspaced was created for the purpose of creating electronic documents, giving its development a suitable reason to support screen reader users. Our tool, \lintex, utilises \lualatex's accessibility support, with emphasis placed on screen reader users.

\subsection{Creating LinTeX}
        - intended purpose
        - (this section has the finished overview diagram)
\subsection{aims and objectives}
\subsection{Limitations of the project overall}
\subsection{research questions/hypothesis}


This is a concise and clear overview of your dissertation(more or less 2-3 pages). You can start off the project description provided of the project that was allocated to you and flesh it out.

Include (1) the problem you have tackled, (2) why this problem is worth addressing, (3) what you did to address it - in broad terms. Detail will come later.



i.e., issue(s) on which the research will focus, shall be clearly identified and described. You shall refer to past research work relevant to the topic and objectives, i.e, of the study. You shall outline where applicable the potential research output with respect to research transfer and uptake by the community.

The introduction says: this is an overview of the project. This is why I did it (the problem) and how I tackled it. It is the \emph{runway} into the project. It lets the marker know what to expect of your report.

If you're doing a research project, this would be the place to include the research questions you plan to address. For example:

\begin{description}
\item [RQ1:] Where did James Bond come from?

\item [RQ2:] Why are pumpkins orange?
\end{description}

If you're doing a project of type 1, include a list of objectives. For example:
\begin{description}
\item [Objective 1:] Provide software to allow James Bond to become invisible.

\item [Objective 2:] Provide software to keep track of all loyalty points in one place.
\end{description}

Provide a `map' of your dissertation. For example: Section \ref{backg} reviews the background literature that was reviewed to inform this project. Then Section \ref{process}.....

\subsection{from nov progress}
There is a lack of accessibility with \latex, largely due to usage of outdated PDFLaTeX as opposed to LuaLaTeX. Usage of this outdated TEX engine necessitates workarounds for desired features and accessibility in resulting PDF documents.

\begin{quote}
What can be done to make \latex documents more accessible?
\end{quote}
LinTeX is a developing tool dedicated to combat this problem by focusing on assisting screen readers' abilities to read the PDF documents produced from \latex documents. It aims to cover this area through suggesting structured tagging throughout a document, suggesting components where it may be lacking. Additionally, it will request users to input appropriate alternative text ('alt text') for non-text components of the document, including figures, images, and tables. Further focus may be placed on features such as ensuring appropriate contrast in any colours in the document, appropriate Sans Serif fonts, and flexible document design, largely depending on what development time allows for. 

Additionally, as a tool focused on promoting accessibility, it too should be accessible. Therefore, the tool will ensure strict adherence to the Web Content Accessibility Guidelines 2 (WCAG 2)%~\cite{wcag} requirements with a focus placed on keeping it a simple and intuitive tool for users to make use of.

\begin{quote}
Why should I use LinTeX?
\end{quote}
While the tool is not limited strictly to the field, it is aimed to assist the many academics that find difficulties in accessing reports and academic literatures. By utilising this tool, you can ensure a wider demographic of users can view your documents, without necessitating an especially large amount of research in accessibility; it has been done for you.